\section{Problems, networks and protocol}\label{sec:methods}

\subsection{Problems and data}\label{sec:methods:data}

Two families of linear elastostatic problems were used, one in two dimensions and one in
three. Both are clamped on the face $x=0$ and carry four load cases that share one
stiffness matrix: a downward and an axial traction on the face opposite the clamp, a shear
traction on the top face, and a downward body force. The magnitude of each load case is
drawn independently as $0.05$ times a number uniform on $[0.5,1.5]$. Young's modulus is
$E=1$ and Poisson's ratio is uniform on $[0.25,0.38]$ per instance; all quantities are
dimensionless. Meshes were generated with Gmsh \cite{Geuzaine2009}, and reference solutions
were computed by sparse direct factorisation.

\paragraph{Two dimensions}
Rectangular plates $[0,l_x]\times[0,l_y]$ in plane stress, with $l_x$ uniform on $[1.5,3]$,
$l_y$ uniform on $[0.8,1.5]$ and zero to three non-overlapping circular holes of radius
between $0.06$ and $0.16$ times $\min(l_x,l_y)$, are meshed with linear triangles of target
size uniform on $[0.05,0.12]$ and assembled with scikit-fem \cite{Gustafsson2020}. The
corpus has 30,000 instances: 256 form the validation set and the first 1,024 of the
remainder the training pool. Two further sets of 2,000 instances, each instance meshed
twice, at the original size and 1.8 or 2.5 times more coarsely, test cross-resolution
transfer (\cref{sec:results2d:other}).

\paragraph{Three dimensions}
Boxes $[0,l_x]\times[0,l_y]\times[0,l_z]$, with $l_x$ uniform on $[1.5,3]$, $l_y$ on
$[0.8,1.5]$, $l_z$ on $[0.6,1.2]$ and zero to three non-overlapping spherical cavities of
radius between $0.08$ and $0.16$ times $\min(l_x,l_y,l_z)$, are meshed with linear
tetrahedra at a characteristic length $l_c$ drawn uniformly from $[0.0579,0.0906]$ per
instance; we call this range in-band. The corpus has 2,000 instances: 256 form the
validation set, the first 1,024 of the remainder the training pool, and the other 720 were
never used. The validation meshes have between \numNodesInMin{} and \numNodesInMax{} nodes
(median \numNodesInMedian{}), with three degrees of freedom per node. A finer set of 320
instances from the same distribution is meshed at $l_c=0.0374$; its first 256 instances,
with between \numNodesFineMin{} and \numNodesFineMax{} nodes (median
\numNodesFineMedian{}), form the evaluation set of \cref{sec:transfer}, and the last 64 are
reserved for few-shot training. No network saw a fine mesh in its in-band training.

\subsection{Networks}\label{sec:methods:networks}

\paragraph{Transformer}
Every mesh node is a token. Its input features are the node's coordinates, centred and
scaled to unit root-mean-square radius; its Dirichlet flags; and its consistent nodal load
for the load case, divided by $F_{\max}$, the largest absolute nodal load component over the
instance's four load cases. Two global descriptors are appended to every token: a load
summary (the total load magnitude and the net load per component, each divided by the number
of nodes and by $F_{\max}$, and the fraction of loaded nodes) and a geometry descriptor
(normalised extents, Poisson's ratio, the number of holes or cavities and their area or
volume fraction, and in two dimensions their mean radius). An encoder of 8 pre-normalised
transformer blocks \cite{Vaswani2017} of width 256 with 8 attention heads applies full
self-attention over the nodes; a decoder maps each node's latent vector, concatenated with
the mean latent vector over the nodes, to the node's displacement through a three-layer
perceptron, and the constrained components are set to zero. Each load case is a separate
sequence; the four sequences of an instance are processed as one batch. The encoder and
decoder have \numParamsTransformer{} million parameters; the implementation also holds two
auxiliary heads, \numParamsAux{} million parameters, which only the secondary objectives of
criteria K3, K4 and K6 use (\ref{app:prereg}).
The decoded field is multiplied by $F_{\max}$, so that the prediction is proportional to the
load amplitude. The two-dimensional runs of July 2026 (\cref{sec:results2d:july},
\ref{app:prereg}) used an earlier version of the code that did not multiply $F_{\max}$ back;
since the load features are normalised by $F_{\max}$, no network in those runs could know the
absolute amplitude of an instance's loads. The runs of September and October 2026 used the
later code.

\paragraph{Graph network}
As a supervised baseline of the most widely used kind we used a MeshGraphNets-style
encode-process-decode network \cite{Pfaff2021} with node and edge encoders, 8
message-passing layers of width 256 over the mesh edges (edge features: relative position
and length), the same node features and the same output scaling as the transformer;
it has \numParamsGraph{} million parameters.

\paragraph{Baselines}
The accuracy tables also report three deterministic baselines: the zero field; a
nearest-neighbour field, the solution of the labelled training instance nearest in the
geometry descriptor, interpolated to the new mesh in normalised coordinates; and a scaled
polynomial, a quadratic ridge regression of the field's direction on the node features with
its norm regressed on the instance's descriptors. We call the last two the naive baselines;
they use the same labelled instances as the supervised networks. The sanity condition of the
two-dimensional gate (\cref{sec:results2d:decisions}) also used a plain polynomial, a
quadratic ridge regression of the displacement itself on the node features.

\subsection{Training}\label{sec:methods:training}

\begin{figure}[tbp]
\centering
\begin{tikzpicture}[
  node distance=5mm and 7mm,
  box/.style={draw, rounded corners=2pt, align=center, font=\footnotesize,
              minimum height=10mm, text width=26mm, inner sep=3pt},
  cost/.style={box, fill=black!8, thick},
  arr/.style={-{Stealth[length=2.2mm]}, semithick}]
\node[box] (i1) {Instance:\\ mesh, material,\\ loads};
\node[box, right=of i1] (a1) {Assemble\\ $K$, $F$};
\node[cost, right=of a1] (s1) {Solve $K\Ustar=F$\\ per load case};
\node[box, right=of s1] (l1) {Loss\\ $\norm{u-\Ustar}_2/\norm{\Ustar}_2$};
\node[box, below=9mm of i1] (i2) {Instance:\\ mesh, material,\\ loads};
\node[box, right=of i2] (a2) {Assemble\\ $K$, $F$};
\node[box, right=27mm of a2] (l2) {Loss\\ $\Pi_h(u)=\tfrac12u^{\top}Ku-F^{\top}u$};
\draw[arr] (i1) -- (a1);
\draw[arr] (a1) -- (s1);
\draw[arr] (s1) -- (l1);
\draw[arr] (i2) -- (a2);
\draw[arr] (a2) -- (l2);
\node[font=\footnotesize, anchor=south west] at ([yshift=1mm]i1.north west) {Supervised};
\node[font=\footnotesize, anchor=south west] at ([yshift=1mm]i2.north west) {Label-free};
\end{tikzpicture}
\caption{The two training routes for the same network. Supervised training needs a
reference solution for every training instance and load case (shaded); label-free
training evaluates the assembled discrete potential energy of the prediction $u$, which by
\cref{lem:exact} differs from the stiffness-norm regression loss by a constant.}
\label{fig:schematic}
\end{figure}

\Cref{fig:schematic} contrasts the two routes. The label-free transformer minimises
\eqref{eq:training} with $L=4$: at each step one training instance is drawn, without
replacement within an epoch, its four load cases are predicted, and the mean of their
energies is the loss; no reference solution is read. The supervised transformer, the same
network, minimises on labelled instances the mean over the load cases of the relative
displacement error,
\begin{equation}
\mathcal{L}_D=\frac{1}{L}\sum_{j=1}^{L}\frac{\norm{u_j-\Ustar_j}_2}{\norm{\Ustar_j}_2},
\label{eq:ld}
\end{equation}
where $u_j$ is the prediction and $\Ustar_j$ the reference solution for load case $j$ of an
instance, and $F_j$ its load vector. In two dimensions the same network was also trained on
the relative stiffness-norm error,
\begin{equation}
\mathcal{L}_K=\frac{1}{L}\sum_{j=1}^{L}\frac{\norm{u_j-\Ustar_j}_K}{\norm{\Ustar_j}_K}
=\frac{1}{L}\sum_{j=1}^{L}g(u_j)^{1/2},
\label{eq:lk}
\end{equation}
which replaces the norm of \eqref{eq:ld} and nothing else; we call this network the
stiffness-norm transformer. Since
$\norm{u_j-\Ustar_j}_K^2=u_j^{\top}Ku_j-2F_j^{\top}u_j+F_j^{\top}\Ustar_j$, the loss
$\mathcal{L}_K$ uses each reference solution only through the number
$\norm{\Ustar_j}_K^2=F_j^{\top}\Ustar_j$, and its gradient with respect to $u_j$,
\begin{equation}
\nabla_{u_j}\mathcal{L}_K=\frac{Ku_j-F_j}{L\,\norm{u_j-\Ustar_j}_K\,\norm{\Ustar_j}_K},
\label{eq:lkgrad}
\end{equation}
is that of the label-free objective of one instance, $(Ku_j-F_j)/L$, weighted per load case.
The supervised transformer with an energy term adds the mean energy of the four load cases
to $\mathcal{L}_D$, with the gradient of the energy scaled to at most the norm of the
gradient of $\mathcal{L}_D$: the update direction is
$\nabla_\theta\mathcal{L}_D+\min\{1,\norm{\nabla_\theta\mathcal{L}_D}/\norm{\nabla_\theta\bar\Pi}\}\nabla_\theta\bar\Pi$,
with $\bar\Pi$ the mean energy. In three dimensions with 16 labelled instances the energy
term was instead $\bar\Pi$ divided by the mean of $\lvert\Pi_h(\Ustar)\rvert$ over the load
cases, with weight one, a policy fixed in the three-dimensional pre-registration after an
exploratory two-dimensional run had shown the scaled form to scatter widely across seeds
with 16 labels (\ref{app:prereg}). The two-dimensional run of July 2026 also trained the
supervised objective from the label-free network's parameters. The graph network is trained
on $\mathcal{L}_D$.

All networks were trained with AdamW \cite{Loshchilov2019} (weight decay $10^{-4}$,
gradient norm clipped at 1) for 200 epochs, with a cosine learning-rate schedule after a
linear warm-up over the first 5\% of the steps, at a peak learning rate of $10^{-3}$
(label-free) or $1.5\times10^{-3}$ (supervised), in single precision with TF32 matrix
products (TensorFloat-32: the inputs of each product rounded to a 10-bit mantissa, the
products accumulated in single precision), each with three seeds (0, 1, 2). The learning rates were fixed before the runs
and not tuned on them; the stiffness-norm transformer used the supervised settings
unchanged. The supervised networks were trained with 16, 64, 256 and 1,024 labelled
instances, the first 16, 64, 256 and 1,024 instances of the training pool. The label-free
network was trained on the same 1,024 instances without their solutions, for 204,800 steps;
its training set is thus the supervised networks' largest one with the labels removed. In
the two-dimensional comparison of \cref{sec:results2d}, the label-free networks are those
trained in the run of 29 September 2026 (\ref{app:e1}) with these settings and the later
code, evaluated without retraining; at a given seed the transformers start from the same
initial weights and, with 1,024 instances, see them in the same sequence of per-epoch
permutations, so that the supervised and the stiffness-norm transformers differ only in the
loss. Except in
the two-dimensional run of July 2026, the graph network was trained with 64 and 1,024
labelled instances only, a saving of machine time fixed in the pre-registrations. The
networks of the secondary checks of \cref{sec:results2d:decisions,sec:results2d:other} were
trained with the settings given where they are reported.

\subsection{Metrics}\label{sec:methods:metrics}

For each validation instance and load case we compute five metrics. The relative
displacement error is $\norm{u-\Ustar}_2/\norm{\Ustar}_2$. The relative energy gap is
$g(u)$ of \eqref{eq:relgap}, computed as $\gap(u)/\lvert\Pi_h(\Ustar)\rvert$. With the
element-wise von Mises stresses $\vm(u)$ and $\vm(\Ustar)$, which are constant on linear
elements, the von Mises error is
$\norm{\vm(u)-\vm(\Ustar)}_2/\norm{\vm(\Ustar)}_2$, the norms taken over the vector of
element values; the peak stress error is
$\lvert\max\vm(u)-\max\vm(\Ustar)\rvert/\max\vm(\Ustar)$; and the critical-region recall is
the share of the $n_c$ elements with the largest $\vm(\Ustar)$ that are also among the
$n_c$ elements with the largest $\vm(u)$, with $n_c$ the nearest integer to a tenth of the
number of elements. Each metric is averaged over the four load cases of an instance and
then over the 256 validation instances; tables give the mean and the sample standard
deviation over the three seeds. The relative energy gap is the primary metric, as
pre-registered: by \cref{prop:stress} it is the strain energy of the stress error relative
to that of the solution, and by \cref{cor:zero} a prediction with $g>1$ is worse than the
zero field.

\subsection{Pre-registration}\label{sec:methods:prereg}

Each run whose results decide a claim of this paper was specified, before it was executed,
in a pre-registration document that fixes the configuration, the metrics, the decision
gates and the kill conditions, which withdraw a claim when they are triggered; both
outcomes of every such run were declared publishable in advance. The configuration is
identified by its SHA-256 hash, and the code refused to start a run whose configuration did
not match the registered hash. Each run was carried out once. The documents of the
three-dimensional runs, of the runs of \ref{app:e1} and \cref{sec:cost:training} and of
the two-dimensional comparison of \cref{sec:results2d} were committed and pushed to the
project's public repository before their runs; tags mark the commits. Those of the
two-dimensional runs of July 2026 carried their configuration hashes, which the code checked
at start-up, but were committed to the repository only on 3 August 2026, after the runs;
their content at the time of the runs rests on the authors' archives, whose hashes the
repository's provenance note records (\ref{app:prereg}). The initial
three-dimensional run needed eight attempts because of memory and process faults; each
restart reused the networks already trained under the same registered configuration, and
the engineering changes the restarts needed are recorded as leaving every value unchanged.
Results that were not pre-registered, an exploratory replication (\cref{sec:results2d:other})
and readings taken after the runs (\cref{sec:results2d:spectrum,sec:results3d:reverse,%
sec:results3d:detection,sec:transfer:amplitude}), are labelled as such.
\Cref{tab:criteria} in \ref{app:prereg} lists every pre-registered criterion of these runs
with its outcome, met or not, under the labels of the pre-registration documents, which
the table explains. An independent script recomputed the three-dimensional verdict from the
report.

Four departures from the plans bear on the reading of the results and are described in
\ref{app:prereg}. First, the initial three-dimensional run (21 September 2026) used a
label-free objective that omitted the output scaling by $F_{\max}$, which inference applied,
so that the trained label-free network returned about \numDfourteenAlpha{} times the
solution (deviation D14). A pre-registered amendment repeated only what depended on that
objective, with the corrected code, and reused the supervised networks, which it had not affected, after
checking their hashes; \cref{sec:results3d,sec:transfer} report the amended run of
28 September 2026. Second, the same amendment, decided after the results of the initial run
were known, assessed one sanity condition only at label budgets of 64 and above; we report
the gate under both readings. Third, the three-dimensional pre-registration text stated
that the label-free network would train on 1,744 instances, whereas the registered
configuration, which is what ran, trained it on the same 1,024 instances as the supervised
networks (deviation D16); the comparison is therefore the one described above. Fourth, a
pre-registered two-dimensional comparison of 31 July 2026 kept, under its decision rule, a
self-supervised term in the pretraining objective; this paper reports the energy objective
without it, a scope chosen after that result was known.

\paragraph{Field plots}
The instances shown in the field plots of \cref{sec:results2d,sec:results3d,sec:transfer}
were chosen by rules fixed in code before any field was drawn: in two dimensions, the
validation instance with the median relative energy gap of the supervised transformer; in
three dimensions, the validation instance with the largest relative energy gap of the
supervised transformer (the same instance in all three seeds), the validation instance with
the median relative energy gap of the label-free transformer, and the fine-mesh instance
with the median displacement error of the label-free transformer; each in seed 0, the median
of $n$ values being the lower one, of rank $\lfloor(n-1)/2\rfloor$ counted from zero in
ascending order. The first three-dimensional rule, and the two-dimensional rule, which was
written after the verdict of the run of 9 October 2026, were chosen after the per-instance
records had been examined, but before any field had been plotted. Each plot shows one load case of its instance, chosen by
the same rule applied to the instance's four load cases (the largest value, or the second
smallest); this was fixed when the plots were made, after the per-load-case values had been
examined but before any field had been drawn. The three-dimensional plots show the three
faces of the box seen in an oblique projection, the free end on the right; the clamped face
and the cavities are hidden.
